
\section{Evaluation and Results}
\subsection{Testing Methodology}
% Description of how the system was tested.

\subsection{Real-world application}
Our system was designed to address challenges with dam systems management, by providing extra analytics to support decisions and prevent failures.

A primary concern in dam operations is the risk of overflow. To mitigate this the system uses multiple sensors, such as pluviometers to gain an accurate insight into the exact quantity of rain, compared to weather estimations, the area is receiving.
Inflows from rivers can also be monitored through the use of a water flow sensor placed at the dam's main water entrance or if the only inflow source is another dam, the system can use flow measurements from the upstream's dam gate as an approximate value.
This gives us valuable insight and allows the model to have an improved performance based on the quantity and sources of data and provides dam technicians with real-time, comprehensive insights into reservoir status.

For optimizing energy production, the system requires two types of sensors and one actuator: a water flow sensor to monitor outflow, a dynamo equipped with a sensor to measure generated electricity, and a flow gate actuator. 
This configuration enables remote and automated control of energy production, allowing the system to adjust operations dynamically in response to changing conditions or to be remotelly actuated by dam technicians.

This would allow our system to provide analytics for technicians to take informed decisions and a way to act on this information in an entirely remote manner helping fix another common issue in dams, their remote or hard to access locations.

\subsection{Result Analysis}
% Detailed breakdown of performance and findings.